\chapter{Base de devises et carry de change}\label{ch:s2:cross-currency-basis-and-fx-carry}

Un panier long sur les trois devises du G10 aux taux d'intérêt les plus élevés face au dollar et short sur les trois plus faibles a rapporté 2.7~\% par an de mai
2002 à décembre 2025, construit ici à partir de taux de change publics et de taux interbancaires. Son ratio de Sharpe était de 0.37, et il a perdu 19.4~\% de septembre à
novembre 2008. Jurek a couvert de tels paniers avec des options hors de la monnaie et a trouvé que les primes de risque de krach représentent au plus un tiers du
rendement du carry trade. L'autre opération de change de ce chapitre est la base de devises. Du, Tepper et Verdelhan ont constaté qu'emprunter des dollars par des swaps de change
coûte plus cher que ne le prévoit la parité des taux d'intérêt couverte, de façon persistante, et surtout pour les contrats inscrits au bilan des banques en fin de trimestre. Sur le marché synthétique de ce chapitre,
un panier de carry couvert par une option de vente abandonne 40~\% de son rendement pour se débarrasser de son skew négatif, et prêter des dollars par le swap rapporte 11.2
points de base par an après un \omterm{def:s2:swap-spread-and-asset-swap-trades:bscost}{coût de bilan}. La construction est \texttt{firm.fxcarry}.

\section{Carry de change~: construction et risque de krach}

\begin{definition}[Panier de carry de change]\label{def:s2:cross-currency-basis-and-fx-carry:basket}
Un \emph{panier de carry de change}\index{panier de carry de change} est long des devises aux taux d'intérêt à court terme les plus élevés et short de celles aux taux les plus bas,
en général à pondérations égales et face à une devise de base commune, rééquilibré à mesure que les taux évoluent~; son rendement est le différentiel de taux perçu auquel s'ajoute la variation
des taux de change.
\end{definition}

Le livre~8 (chapitre~20) a construit la stratégie de portage (carry) sur l'ensemble des classes d'actifs et sur ses krachs. Ici, la version sur devises est construite à partir de données publiques. Chaque mois,
sur les taux interbancaires à trois mois du mois précédent, le panier classe neuf devises (euro, yen, livre sterling, franc suisse, dollars australien et néo-zélandais,
dollar canadien, couronnes suédoise et norvégienne) selon leur taux moins le taux américain. Il est long des trois premières et short des trois dernières face au
dollar. De mai 2002 à décembre 2025 (284 mois), il a rapporté 2.7~\% par an avec une volatilité de 7.2~\%, un ratio de Sharpe de 0.37 et une skewness de
$-0.60$ (la \cref{fig:s2:cross-currency-basis-and-fx-carry:g10}). Le différentiel de taux a rapporté 3.1~\% par an et les taux de change en ont repris 0.4~\%. Son
pire mois a été octobre 2008, $-10.5~\%$. Le dollar néo-zélandais a été long pendant 271 des 284 mois, le dollar australien pendant 218, la couronne norvégienne pendant 188. Le
franc suisse a été short pendant les 284 mois, le yen pendant 177, la couronne suédoise pendant 166, l'euro pendant 159.

\begin{omfigure}
\begin{tikzpicture}
\begin{axis}[width=11cm,height=5.4cm,xlabel={\small année},ylabel={\small rendement logarithmique cumulé (\%)},tick label style={font=\footnotesize},
  xmin=2002,xmax=2026.5,ymin=0,ymax=65,grid=major,grid style={gray!25},table/col sep=comma,/pgf/number format/1000 sep={}]
\addplot[omThm,thick] table[x=year,y=cum]{figdata/strategies-2/14-cross-currency-basis-and-fx-carry/g10.csv};
\end{axis}
\end{tikzpicture}

\omcaption{Panier de carry du G10 face au dollar (long sur les trois taux interbancaires à trois mois les plus élevés, short sur les trois plus bas), rendement logarithmique cumulé,
de mai 2002 à décembre 2025, à partir des taux de change de FRED et des taux interbancaires de l'OCDE. Statistiques dérivées uniquement. Données~:
\texttt{s2\_fetch\_g10}.}\label{fig:s2:cross-currency-basis-and-fx-carry:g10}
\end{omfigure}

Le panier fondé sur des données publiques n'est qu'approximatif. Les taux de l'OCDE sont mensuels et diffèrent d'un pays à l'autre selon leur mode de mesure. Un trader utiliserait des opérations de change à terme, dont
les points de terme donnent exactement le différentiel de taux, et en paierait les fourchettes. La configuration est celle qu'on connaît~: des gains réguliers, puis un krach lorsque les
devises à haut rendement chutent ensemble.

\subsection*{Couvertures contre les krachs}

\begin{definition}[Carry couvert contre les krachs]\label{def:s2:cross-currency-basis-and-fx-carry:hedged}
Le \emph{carry couvert contre les krachs}\index{carry couvert contre les krachs} est une position de carry combinée à des options hors de la monnaie qui rapportent lorsque les devises à haut rendement
chutent fortement face aux devises de financement, de sorte que le rendement de la position ne subisse plus le krach~; le coût des options mesure le
prix du risque de krach que supporte le trade non couvert.
\end{definition}

Jurek a utilisé les prix d'options sur le G10 de 1990 à 2012. Les trades de carry ont affiché des ratios de Sharpe aussi élevés que les facteurs actions, voire davantage, et la couverture par
options hors de la monnaie a montré que les rendements n'étaient pas un peso problem. Les primes de risque de krach en représentaient au plus un tiers.

\texttt{firm.fxcarry} reprend la même comparaison sur les dix devises synthétiques de l'univers de contrats à terme du livre~8, qui portent une prime et trois
krachs programmés (le \cref{lst:s2:cross-currency-basis-and-fx-carry:hedge}). Le panier a rapporté 1.30~\% par an, avec un ratio de Sharpe de 0.23, une skewness de $-0.39$ et
un pire mois de $-6.6~\%$. Une option de vente mensuelle sur le panier, de prix d'exercice fixé à 1.5 écart-type sous zéro et valorisée à 1.2 fois la volatilité glissante, a coûté
1.14~\% par an et a rapporté 0.63~\%. Couvert, le panier a rapporté 0.78~\% par an, un ratio de Sharpe de 0.15, sans skew (0.01) et un pire mois de $-4.3~\%$
(la \cref{fig:s2:cross-currency-basis-and-fx-carry:hedge}). Les krachs programmés ont duré une centaine de jours chacun et ont coûté 14.1~\%, 15.5~\% et 14.1~\% au
panier. Sur les mois qu'ils couvraient, la couverture a ramené le premier de $-12.0~\%$ à $-7.9~\%$ et le troisième de $-14.4~\%$ à $-8.5~\%$. Elle n'a presque rien
fait pour le deuxième ($-14.7~\%$ à $-14.3~\%$), qui a reculé un peu chaque mois, sans jamais descendre assez sous le prix d'exercice. Le coût net de la couverture a représenté 40~\% du
rendement de carry, plus que le tiers de Jurek. C'est le résultat d'un choix délibéré sur les prix d'options.

\begin{omfigure}
\begin{tikzpicture}
\begin{axis}[width=11cm,height=5.4cm,xlabel={\small année},ylabel={\small rendement logarithmique cumulé (\%)},tick label style={font=\footnotesize},
  xmin=0,xmax=30,ymin=-15,ymax=50,grid=major,grid style={gray!25},table/col sep=comma,
  legend style={font=\scriptsize,at={(0.03,0.97)},anchor=north west}]
\addplot[omThm,thick] table[x=year,y=raw]{figdata/strategies-2/14-cross-currency-basis-and-fx-carry/hedge.csv};
\addplot[omDef,thick] table[x=year,y=hedged]{figdata/strategies-2/14-cross-currency-basis-and-fx-carry/hedge.csv};
\legend{panier de carry,avec puts mensuels}
\end{axis}
\end{tikzpicture}

\omcaption{Le panier de carry synthétique sur devises, sur trente ans, avec trois krachs programmés, non couvert et avec une option de vente mensuelle dont le prix d'exercice est fixé à 1.5
écarts-types en dessous. Données~: \texttt{s2\_fxcarry.cumulative}.}\label{fig:s2:cross-currency-basis-and-fx-carry:hedge}
\end{omfigure}

\section{La base de devises}

La parité des taux d'intérêt couverte (livre~2, chapitre~16) veut que s'endetter directement en dollars et emprunter une autre devise puis la convertir en dollars par un swap coûtent
le même prix. La base de devises mesure l'écart entre les deux~; une base négative signifie que la voie du swap coûte davantage, de sorte que les dollars sont chers
via les swaps de change. Du, Tepper et Verdelhan ont constaté des écarts importants, persistants, systématiques et non expliqués par le risque de crédit ni par les coûts de transaction. Ils
étaient corrélés à d'autres spreads obligataires.

\begin{definition}[Trade de financement par la base]\label{def:s2:cross-currency-basis-and-fx-carry:funding}
Un \emph{trade de financement par la base}\index{trade de financement par la base} prête des dollars par des swaps de change (en recevant l'autre devise et en la restituant à terme) lorsque
la base de devises rend cette voie plus rémunératrice que le prêt direct de dollars. Le gain vaut moins la base, diminué du \omterm{def:s2:swap-spread-and-asset-swap-trades:bscost}{coût de bilan} que
la position utilise.
\end{definition}

La base synthétique se situe à $-20$ points de base par an, avec du bruit, auxquels s'ajoutent des fenêtres de fin de trimestre~: dans les dix derniers jours de négociation de chaque trimestre, elle est
30 points de base de plus négative, 60 en fin d'année. Elle s'est établie en moyenne à $-26.2$ points de base sur dix ans, $-49.8$ dans les fenêtres de fin de trimestre et $-80.4$ dans les
celles de fin d'année. Un prêteur dont le \omterm{def:s2:swap-spread-and-asset-swap-trades:bscost}{coût de bilan} atteint 15 points de base par an a gagné 11.2
points de base par an en prêtant en permanence. Prêter seulement pendant les fenêtres, soit 40 jours par an, a rapporté 6.7, à un taux de 42.5 points de base par an lorsque les dollars étaient prêtés.

\section{Fins de trimestre}

La configuration de fin de trimestre est l'argument du \omterm{def:s2:swap-spread-and-asset-swap-trades:bscost}{coût de bilan} sous sa forme la plus simple. Rien ne change côté crédit ni côté taux le 31 mars, mais les banques qui publient
leur bilan à cette date le réduisent sur la période, et les dollars qu'elles auraient prêtés par le swap deviennent rares. Les fenêtres du chapitre sont
programmées. Les fenêtres réelles n'ont que l'ampleur des règles de bilan qui les provoquent. Un prêteur qui n'est pas lui-même mesuré à ces dates est rémunéré pour
fournir le bilan que d'autres retirent.

\section{Trades guidés par le financement}

Les deux moitiés du chapitre se rejoignent dans le financement. Un trade de carry se finance dans la devise à faible taux. Lorsque les dollars sont chers par les swaps, les points
de terme qui fixent le carry intègrent la base. Une position de carry couverte sur le marché des swaps paie la base, et une position financée en cash ne la paie pas. Les portefeuilles de carry
s'effondrent aussi lorsque le financement se resserre. Les devises à haut rendement chutent quand leurs détenteurs vendent pour rembourser les prêts libellés dans la devise de financement. Un filtre de momentum qui coupe le
panier quand sa tendance propre devient négative est la parade habituelle. Il abandonne une partie du carry dans les phases calmes et manque les premiers jours de tout krach.

\section{Fiches de stratégie}

\begin{strategyfile}{Panier de carry G10}\label{strat:s2:cross-currency-basis-and-fx-carry:basket}
\sfield{Qui rémunère, et pourquoi} Des investisseurs qui paient pour détenir des devises sûres à faible rendement~; le risque d'un effondrement des devises à haut rendement lors d'une fuite vers la sécurité.
\sfield{Instruments et plateformes} Contrats à terme de change~; contrats à terme.
\sfield{Signal} Points de terme (l'écart de taux).
\sfield{Dimensionnement et exécution} Trois devises en tête à l'achat, trois en queue à la vente, à pondération égale~; mensuel.
\sfield{Coûts} Spreads à terme~; roll.
\sfield{Scénario d'échec} Krachs où les devises à haut rendement chutent ensemble~: $-19.4~\%$ à l'automne 2008 pour le panier fondé sur des données publiques.
\sfield{Horizon, capacité, infrastructure} Des mois à plusieurs années~; grande capacité.
\sfield{Backtest honnête} Points de terme, et non taux interbancaires~; coûts de transaction~; le krach de 2008 dans l'échantillon.
\sfield{Sources} Jurek (2014)~; le panier FRED/OECD de ce chapitre~: 2.7~\% par an, ratio de Sharpe 0.37, 2002--2025.
\end{strategyfile}

\begin{strategyfile}{Carry couvert contre les krachs avec options}\label{strat:s2:cross-currency-basis-and-fx-carry:hedged}
\sfield{Qui rémunère, et pourquoi} Comme pour le carry, moins la prime de krach versée aux vendeurs d'options.
\sfield{Instruments et plateformes} Positions de carry et options de change hors de la monnaie sur chaque paire ou sur le panier.
\sfield{Signal} Le carry~; la couverture est programmée.
\sfield{Dimensionnement et exécution} Options mensuelles dont le prix d'exercice se situe à un ou deux écarts-types.
\sfield{Coûts} Primes d'options et fourchettes.
\sfield{Scénario d'échec} Krachs lents qui n'atteignent jamais le prix d'exercice.
\sfield{Horizon, capacité, infrastructure} Mensuel~; accès aux options.
\sfield{Backtest honnête} Prix des options issus du smile de volatilité de l'époque, et non d'une volatilité plate.
\sfield{Sources} Jurek (2014)~: les primes de krach atteignent au plus un tiers des rendements du carry~; ce chapitre~: la couverture coûte 40~\% et élimine le skew.
\end{strategyfile}

\begin{strategyfile}{Basis trade de fin de trimestre}\label{strat:s2:cross-currency-basis-and-fx-carry:qe}
\sfield{Qui rémunère, et pourquoi} Des banques qui réduisent leur bilan autour des dates de publication.
\sfield{Instruments et plateformes} Swaps de change à cheval sur la fin de trimestre.
\sfield{Signal} Le calendrier~; la base cotée pour le passage de fin de période.
\sfield{Dimensionnement et exécution} Prêter des dollars à travers le passage de fin de période~; limité par son propre bilan à cette date.
\sfield{Coûts} \omterm{def:s2:swap-spread-and-asset-swap-trades:bscost}{Coût de bilan} propre au passage de fin de période.
\sfield{Scénario d'échec} Règles de publication qui moyennent sur le trimestre.
\sfield{Horizon, capacité, infrastructure} Quelques jours à quelques semaines.
\sfield{Backtest honnête} Prix des swaps aux dates exactes du passage de fin de période.
\sfield{Sources} Du, Tepper et Verdelhan (2018)~; ce chapitre~: 42.5 points de base par an pendant les prêts sur les fenêtres.
\end{strategyfile}

\begin{strategyfile}{Arbitrage de funding sur la base}\label{strat:s2:cross-currency-basis-and-fx-carry:funding}
\sfield{Qui rémunère, et pourquoi} Des emprunteurs de dollars par swaps de change (banques non américaines et investisseurs couvrant des actifs en dollars) qui paient au-dessus de la parité.
\sfield{Instruments et plateformes} Swaps de change et swaps de devises~; financement en dollars.
\sfield{Signal} La base comparée à son propre financement et au \omterm{def:s2:swap-spread-and-asset-swap-trades:bscost}{coût de bilan}.
\sfield{Dimensionnement et exécution} Selon la capacité du bilan~; échéances appariées.
\sfield{Coûts} \omterm{def:s2:swap-spread-and-asset-swap-trades:bscost}{Coût de bilan}~; lignes de contrepartie.
\sfield{Scénario d'échec} Une tension sur le financement du prêteur lui-même.
\sfield{Horizon, capacité, infrastructure} Des mois~; financement au niveau de la trésorerie.
\sfield{Backtest honnête} Imputer le \omterm{def:s2:swap-spread-and-asset-swap-trades:bscost}{coût de bilan} réel du prêteur.
\sfield{Sources} Du, Tepper et Verdelhan (2018)~; ce chapitre~: 11.2 points de base par an après un coût de 15 points de base.
\end{strategyfile}

\begin{strategyfile}{Carry avec filtre de momentum}\label{strat:s2:cross-currency-basis-and-fx-carry:momentum}
\sfield{Qui rémunère, et pourquoi} Comme pour le carry, avec une certaine protection contre les krachs, payée par le carry abandonné.
\sfield{Instruments et plateformes} Comme pour le carry.
\sfield{Signal} Le carry, coupé lorsque le rendement glissant du panier est négatif.
\sfield{Dimensionnement et exécution} Réduire ou clôturer lorsque le filtre se déclenche.
\sfield{Coûts} Un turnover plus élevé.
\sfield{Scénario d'échec} Les faux signaux répétés~; krachs rapides qui débutent avant le déclenchement du filtre.
\sfield{Horizon, capacité, infrastructure} Des mois.
\sfield{Backtest honnête} Paramètres du filtre choisis avant le test.
\sfield{Sources} Aucun chiffre de performance n'a été vérifié.
\end{strategyfile}

\section{Tutoriel~: cotée en fin de trimestre}\label{tut:s2:cross-currency-basis-and-fx-carry:tut}

\begin{tutorial}
\textbf{Objectif.} Construire des paniers de carry, couvrir leurs krachs par des options, simuler une base de devises avec des fenêtres de fin de trimestre et prêter dans cette base~;
construire le panier G10 à partir de données publiques. \textbf{État final~:} les chiffres du texte et les deux figures.
\begin{tutsteps}
\item \textbf{Panier et couverture}.
\omcode{code/firm/fxcarry/firm_fxcarry.py}{27}{64}{Paniers de carry et un put mensuel sur le panier.}\label{lst:s2:cross-currency-basis-and-fx-carry:hedge}
\item \textbf{Résultats du carry}.
\omcode{code/strategies-2/14-cross-currency-basis-and-fx-carry/python/s2_fxcarry.py}{45}{55}{Statistiques sans et avec couverture, et mois de krach.}\label{lst:s2:cross-currency-basis-and-fx-carry:results}
\item \textbf{Exécuter} \texttt{s2\_fetch\_g10.py} une fois, puis \texttt{carry\_results()}, \texttt{basis\_results()} et \texttt{fig\_fxcarry.py}.
\end{tutsteps}
\textbf{Pour aller plus loin.} Couvrir chaque paire par sa propre option au lieu du panier~; ajouter un filtre de momentum et comparer les krachs~; lier la base
au financement du panier de carry.
\end{tutorial}

\section{Construction~: carry de change et base}\label{bld:s2:cross-currency-basis-and-fx-carry:fxcarry}

\begin{build}
\bfield{Rôle} Paniers de carry, carry couvert par options, base de devises avec fenêtres de fin de trimestre, et \omterm{def:s2:cross-currency-basis-and-fx-carry:funding}{trades de financement par la base}.
\bfield{Interface} \texttt{carry\_basket(r, carry, k, every)}, \texttt{crash\_hedge(basket, every, strike\_sd, vol\_mult, lookback)},
\texttt{BasisConfig(\dots)}, \texttt{simulate\_basis(cfg)}, \texttt{basis\_trade(b, cfg, days)}.
\bfield{Règles} Rééquilibrer sur le carry de la veille~; puts évalués au début de chaque période sur la volatilité glissante~; les basis trades rapportent l'opposé de la base moins le
\omterm{def:s2:swap-spread-and-asset-swap-trades:bscost}{coût de bilan}.
\bfield{Tests d'acceptation} \texttt{code/firm/fxcarry/tests/}~: les positions du panier à la main~; le prix et le paiement du put à la main~; les fenêtres de base
et le couru quotidien du trade.
\bfield{Pour aller plus loin} Options par paire~; points de terme avec la base~; krachs de financement.
\end{build}

\begin{omsources}
\item J.~W.~Jurek, ``Crash-neutral currency carry trades'', \emph{Journal of Financial Economics} 113(3), 2014.
\item W.~Du, A.~Tepper et A.~Verdelhan, ``Deviations from covered interest rate parity'', \emph{Journal of Finance} 73(3), 2018.
\item Board of Governors H.10 exchange rates and OECD interbank rates, via FRED.
\end{omsources}

\section{Exercices}

\begin{exercise}[$\star$]\label{exo:s2:cross-currency-basis-and-fx-carry:1}
Le panier a rapporté 2.7~\% par an avec une volatilité de 7.2~\%. Calculer son ratio de Sharpe, sans taux sans risque.
\end{exercise}

\begin{exercise}[$\star$]\label{exo:s2:cross-currency-basis-and-fx-carry:2}
La base vaut $-50$ points de base par an et le \omterm{def:s2:swap-spread-and-asset-swap-trades:bscost}{coût de bilan} d'un prêteur est de 15. Calculer ce que rapporte un prêt de dix jours ouvrés, en points de base du
montant.
\end{exercise}

\begin{exercise}[$\star$]\label{exo:s2:cross-currency-basis-and-fx-carry:3}
Une couverture coûte 1.14~\% par an et verse 0.63~\%. Calculer son coût net en proportion d'un rendement du carry de 1.30~\%.
\end{exercise}

\begin{exercise}[$\star\star$]\label{exo:s2:cross-currency-basis-and-fx-carry:4}
Pourquoi le franc suisse était-il en position courte dans chaque mois du panier de données publiques~?
\end{exercise}

\begin{exercise}[$\star\star$]\label{exo:s2:cross-currency-basis-and-fx-carry:5}
Pourquoi le put mensuel n'a-t-il pas protégé contre le second krach synthétique~?
\end{exercise}

\begin{exercise}[$\star\star$]\label{exo:s2:cross-currency-basis-and-fx-carry:6}
Pourquoi une base de devises négative persiste-t-elle~?
\end{exercise}

\begin{exercise}[$\star\star\star$]\label{exo:s2:cross-currency-basis-and-fx-carry:7}
\emph{Programmation.} Exécuter \texttt{carry\_results(vol\_mult=1.0)}, en évaluant les puts à la volatilité glissante, sans prime. Que coûte maintenant la couverture, et
que dit cela de l'endroit où se situe la prime de krach~?
\end{exercise}

\begin{exercise}[$\star\star\star$]\label{exo:s2:cross-currency-basis-and-fx-carry:8}
\emph{Trouver la faille.} «~Le carry a rapporté 2.7~\% par an avec un ratio de Sharpe de 0.37 sur 23 ans~; nous le porterons à une volatilité de 20~\% pour 7.5~\% par an.~»
\end{exercise}

\section{Problème~: Valorisation à la fin du trimestre}

\begin{problem}[{Problème du week-end --- devises et bilans}]\label{pb:s2:cross-currency-basis-and-fx-carry:1}
Les marchés synthétiques du chapitre, le panier de données publiques et le registre public.

\medskip\noindent\textbf{Partie I --- Carry.}
\begin{enumerate}
\item Définir un \omterm{def:s2:cross-currency-basis-and-fx-carry:basket}{panier de carry de change} et le construire à partir de données publiques.
\item Donner son rendement, son risque et son krach.
\item Quelles devises le panier détenait-il, et à quelle fréquence~?
\item Que Jurek a-t-il trouvé~?
\end{enumerate}

\medskip\noindent\textbf{Partie II --- Couvertures.}
\begin{enumerate}[resume]
\item Définir le \omterm{def:s2:cross-currency-basis-and-fx-carry:hedged}{carry couvert contre les krachs}.
\item Donner les statistiques du panier synthétique avec et sans la couverture.
\item Que fait la couverture lors de chaque krach~?
\item Pourquoi le coût de la couverture mesure-t-il la prime de krach~?
\end{enumerate}

\medskip\noindent\textbf{Partie III --- La base.}
\begin{enumerate}[resume]
\item Définir la parité des taux d'intérêt couverte et la base de devises.
\item Que trouvent Du, Tepper et Verdelhan~?
\item Définir le \omterm{def:s2:cross-currency-basis-and-fx-carry:funding}{trade de financement par la base} et donner son rendement synthétique.
\item Qu'apporte la fenêtre de fin de trimestre~?
\end{enumerate}

\medskip\noindent\textbf{Partie IV --- Le verdict.}
\begin{enumerate}[resume]
\item Énoncer le \emph{résultat nommé}~: le ratio de Sharpe du panier de carry avec et sans couvertures contre les krachs, et le rendement de fin de trimestre du basis trade.
\item Comment le carry et la base se rejoignent-ils dans le financement~?
\item Que changerait un filtre de momentum~?
\item Comment dimensionner un book de carry~?
\item Comment mener un backtest honnête du basis trade~?
\item Quelle fiche de stratégie a besoin de son propre bilan~?
\item Quel est le lien entre ce chapitre et le chapitre~12~?
\item En une phrase~: que valorise la base de devises~?
\end{enumerate}
\end{problem}

\section{Questions d'entretien}

\begin{interviewq}[$\star$ \iqroles{trader}]\label{iq:s2:cross-currency-basis-and-fx-carry:1}
Qu'est-ce que le carry trade de change, et pourquoi fonctionne-t-il~?
\end{interviewq}

\begin{interviewq}[$\star\star$ \iqroles{chercheur}]\label{iq:s2:cross-currency-basis-and-fx-carry:2}
Comment mesureriez-vous la prime de risque de krach dans le carry de change~?
\end{interviewq}

\begin{interviewq}[$\star\star$ \iqroles{trader}]\label{iq:s2:cross-currency-basis-and-fx-carry:3}
La base de devises à trois mois est de $-40$ points de base. Que signifie-t-elle, et qui la paie~?
\end{interviewq}

\begin{interviewq}[$\star\star$ \iqroles{risque}]\label{iq:s2:cross-currency-basis-and-fx-carry:4}
Comment stresseriez-vous un book de carry G10~?
\end{interviewq}

\begin{interviewq}[$\star\star$ \iqroles{développeur}]\label{iq:s2:cross-currency-basis-and-fx-carry:5}
Concevez les données sous-jacentes à un book de carry~: taux, cours à terme, comptant, jours fériés, et ce qui peut mal tourner en fin de mois.
\end{interviewq}

\begin{interviewq}[$\star\star\star$ \iqroles{chercheur}]\label{iq:s2:cross-currency-basis-and-fx-carry:6}
Dérivez le cours de change à terme à partir de la parité des taux d'intérêt couverte, et montrez comment une base $b$ y entre.
\end{interviewq}
